\section{Encoder Implementation and Benchmark Details}
\label{app:implementation}

This appendix supports Section~\ref{sec:implementation}.
It gives the full implementation description, memory accounting,
measurement protocols, and ordinary-encoding comparison.
We evaluate the transposed BCH--IMT encoders of
Section~\ref{sec:implementation} on batches of $N$ input blocks, with
$N=2K$ at rate $1/2$ and $N=4K$ at rate $1/4$.

\paragraph{Shared Implementation and Optimized Kernels.}
The configurations share the routing machinery and an IMT template.
The transpose uses
\[
 V_i=U_i+C^\top r_{i+1},\qquad r_i=A^\top U_i+M_i^\top r_{i+1}.
\]
In particular, $A^\top$ consumes the raw input $U_i$.
The balanced expansion is evaluated by pruned SIMD zeta stages and fixed
nibble tables. Sparse feedback compiles to unrolled XORs; each output
block combines five state blocks for the $(128,19)$ rate-$1/2$ inner,
or three at rate $1/4$. The small-length variant uses its own fixed maps.
The transvection update uses XORs on selected state blocks.
The selected BCH-256 circuit evaluates four rows together with AVX-512,
using a generated $3758$-XOR circuit; the AVX2 fallback uses a
$3029$-XOR circuit on two rows. The rate-$1/4$ constituent has its own generated
transpose circuit. The online call has no allocation or type-erased
callback dispatch; configuration and layout dispatch occur once per encoding.

Packed $24$-bit routing indices and workspace-local page routing reduce
data-movement costs. At $K=2^{20}$ the rate-$1/2$ implementation uses
$2048$-row tiles and the rate-$1/4$ implementation uses $4096$-row tiles.
The smaller rate-$1/2$ measurements use direct routing and four-row SIMD
packing, avoiding the intermediate tile scatter.
These choices were tuned on the evaluation host.

\paragraph{Correctness and Certificate Binding.}
Symbolic checks verify the generated outer matrices and transpose circuits,
the pruned zeta schedule, and both independent inner maps.
Runtime tests compare the complete optimized output with dense outer
multiplication and a separately implemented scalar inner recurrence.
The checks cover routing layouts, tile sizes, setup compaction,
boundary inputs, buffer validation, and agreement between SIMD backends.
Source and binary identities, exact map columns,
and the outer construction bind each reported timing to its corresponding
finite certificate. A timing alone does not establish distance.

\paragraph{Measurement Method.}
Measurements use one Ryzen 9 7950X thread on Linux, pinned to CPU~15, with
GCC~15.2.0 and boost disabled. The rate-$1/2$ SPIN build uses
\texttt{-O3 -DNDEBUG -march=x86-64-v3 -mtune=znver4},
carry-less multiplication instructions, and AVX-512 for the BCH circuit.
The rate-$1/4$ and external comparison builds use \texttt{-march=znver4}.
Each reported time is the median of three independent serial process medians.
Each process performs three warmups and $101$ timed calls. Buffers are initialized
once per process, with no copy or reset between calls.
Setup, allocation, correctness checks, and output hashing
are outside the online interval. These are steady-state encoder latencies,
not cold-cache or full-application timings. All setup is precomputed;
no heuristic refresh is included.

\paragraph{A Small-Length Variant.}
\label{app:small-k-variant}
At $K=2^{16}$, we keep the BCH-256 outer and routing but use
$(t,s)=(64,12)$ with two independent transvections per inner step.
The smaller maps reduce arithmetic cost; the extra transvection improves
state mixing. A separate full first-moment certificate covers all
$512$ possible active-row counts and gives a $49.328$-bit margin at
relative distance $0.10$.
This is an additional finite operating point outside the
one-transvection family of Theorem~\ref{thm:finite-bch-spin}.
The repository documentation describes its fixed maps and
certificate\repositorycite.

\begin{table}[tb]
\centering
\caption{Online time in milliseconds for the batched transposed encoder.
The rate-$1/2$ row uses the small-length variant at $K=2^{16}$;
the other cells use IMT $(t,s)=(128,19)$ with their specified maps.
A dash means that no certificate-matched timing is reported here.}
\label{tab:finite-bch-performance}
\begin{tabular}{@{}lrrr@{}}
\toprule
Outer and rate & $K=2^{16}$ & $K=2^{18}$ & $K=2^{20}$ \\
\midrule
BCH-256, $1/2$ & 0.340 & 1.460 & 9.289 \\
BCH-128, $1/4$ & -- & -- & 15.254 \\
\bottomrule
\end{tabular}
\end{table}

At rate $1/2$ and $K=2^{20}$, the three process medians range from
$9.158$ to $9.321$\,ms. Precomputed setup occupies $12.13$\,MiB,
workspace $40$\,MiB, and input/output buffers a separate $48$\,MiB.
The rate-$1/4$ setup occupies $24.25$\,MiB and workspace $72$\,MiB;
its process medians range from $15.204$ to $15.313$\,ms.
Setup latency is not reported. We report no encoding timings at
$K=2^{22}$ or $2^{24}$.
The repository provides the encoder implementation, its documentation,
and the measurement and certificate data for both rates\repositorycite.

\subsection{Ordinary Encoding}
\label{sec:ordinary-encoding}

Polynomial commitments require the ordinary encoder: $K$ input blocks
expand to $2K$ output blocks. We implement this direction for the same
selected SPIN map and measure it separately. Each block again contains
$128$ parallel binary instances. These measurements include the outer
encoder, routing, and inner recurrence, but no commitment or opening work.

The ordinary encoder reuses the same local circuits and direct or tiled routing
strategy. It gathers the inner input and emits the output in one pass,
retaining the raw input for the independent feedback map. A generated XOR
circuit evaluates that map. Both directions use the same transvection
updates and workspace-local page routing. Dense-oracle checks and the identity
$\langle E(x),q\rangle=\langle x,E^\top(q)\rangle$ check that the
implementation preserves the linear map.

\begin{table}[tb]
\centering
\caption{Ordinary and transposed SPIN encoding, in milliseconds, for
$128$ parallel binary instances. Setup is fully precomputed, input is
preserved, and output is separate; neither direction includes PCS work.}
\label{tab:ordinary-encoding}
\begin{tabular}{@{}lrrr@{}}
\toprule
Direction & $K=2^{16}$ & $K=2^{18}$ & $K=2^{20}$ \\
\midrule
% Generated by build_application_tables.py; data SHA-256 07c6005a5a76085c42de1fb37977c1d8b2633a7508deef3ce9d56b95db534297
% Optimized precomputed runs SHA-256 fe4afcc33fec10d1a802500394991ad6091b03fef646bbd0f4f62f68c540b087
Ordinary & 0.348 & 1.601 & 9.031 \\
Transposed & 0.340 & 1.460 & 9.289 \\
\bottomrule

\end{tabular}
\end{table}

The experiment uses the same Ryzen host, CPU affinity, and GCC version
as above, with boost disabled. Each direction has independent input and
workspace buffers. Separate processes alternate direction order, each
performing three warmups and $101$ timed calls.
Table~\ref{tab:ordinary-encoding} reports medians of three process medians.
Setup, allocation, and output checks are outside the interval.
All runs execute serially.

Table~\ref{tab:ordinary-encoding} shows similar cost for both directions
($9.031$\,ms ordinary at $K=2^{20}$).
The same experiment measures in-place transposed encoding at
$K=2^{18}$ in $1.419$\,ms, versus $1.460$\,ms with separate output.
Dividing $2^{18}$ output blocks by this in-place latency gives approximately
$185$ million $128$-bit correlation blocks per second.
The ordinary encoder uses $256$-row tiles at $K=2^{20}$ and direct routing
at the smaller lengths. Its workspace occupies $33$\,MiB at $K=2^{20}$.
The OT, PCS, and Flock measurements use their own integrated encoder
builds; we do not rescale them using these encoder timings.

\subsection{Comparison with Other Transposed Encoders}
\label{app:transposed-comparison-details}

We compare SPIN with BAA, introduced by Kolesnikov et al.\ at CRYPTO
2026~\cite{KolesnikovEtAl2026BA}, and the newer chosen-block
instantiations~\cite{cryptoeprint:2026/1903}. Chosen-block BAA revises
BAA, including by reusing a spectrum-controlled outer constituent. Table~\ref{tab:transposed-comparison}
includes both BAA variants. The other baselines are
Expand--Convolute~\cite{RaghuramanRindalTanguy2023EC} and binary
repeat--accumulate--accumulate (RAA)~\cite{blaze25}.

Every call evaluates a binary
transposed encoder on $128$ parallel instances. No weighted extension-field
RAA or ordinary-encoder timing is included.

The chosen-block BAA rows use the optimized Golay $[24,12,8]$ and RM$(2,5)$ kernels
from the chosen-block BAA implementation. Expand--Convolute uses the
nonsystematic implementation with the conservative profile, expander weight $w=33$ and memory $m=25$.
The implementation takes $24$ random taps plus one fixed tap and applies
one convolution pass. Its native interface modifies the input buffer and
writes the result to a preallocated output buffer. Our binary RAA adapter
uses two accumulator/permutation layers and a repetition transpose.
Chosen-block BAA and RAA use four-round Feistel permutations, with cycle walking for
non-power-of-two domains. SPIN uses its packed routing tables.
The measured times include each implementation's routing costs.

\paragraph{Common Timing Protocol.}
Except for the original-BAA reference, all rows use GCC~15.2.0, with the
per-implementation flags above, on the same Ryzen host and CPU. Each cell uses three serial processes, each with
three warmups and $101$ timed calls. Each process allocates and initializes
its buffers once, then repeatedly invokes the native encoder interface.
There is no input copy or reset between calls. Expand--Convolute operates
on the buffer left by the previous call; the other interfaces preserve
their input. Code setup, caller buffer allocation, and output hashing
are outside the interval; all work inside each encoder call is timed.
The chosen-block BAA message lengths are $2^m-64$ for $m\in\{16,18,20\}$, allowing
both constituents to tile the same input. All other rows use $K=2^m$.
RAA processes $3K$ or $4K$ input blocks, versus $2K$ for the rate-$1/2$ rows.

The combined results appear in the main comparison,
Table~\ref{tab:transposed-comparison}. Its RAA footnote states the different
reference length and the optional setup-testing guarantee.

At the largest measured length, rate-$1/2$ SPIN with IMT takes $9.289$\,ms.
The original rate-$1/2$ BAA implementation takes approximately $32$\,ms,
giving the approximately $3.4\times$ speedup stated in the abstract.
The newer chosen-block implementations improve on that baseline:
chosen Golay and RM BAA take $2.58$ and $2.98$ times as long as SPIN,
respectively. The conservative EC profile takes $181.825$\,ms.
The repetition-$3$ permutation domain requires cycle walking; the
repetition-$4$ domain does not. The faster repetition-$4$ timings therefore
do not show that repetition-$4$ RAA is intrinsically faster.

\paragraph{Distance Context.}
At the largest length, rate-$1/2$ SPIN's certified margin is $50.062$ bits at
relative distance $0.10$. The chosen-block BAA evaluation reports full
first-moment margins of $46.386$ and $47.093$ bits for Golay and RM,
respectively, at the same distance target and $K=2^{20}-64$.
Original BAA's conservative $\sigma=32$ profile has relative distance
$0.10$ with a $20$-bit margin at $K=2^{20}$~\cite[Fig.~10]{KolesnikovEtAl2026BA}.
For EC, we use the conservative profile reported by Kolesnikov
et al.~\cite{KolesnikovEtAl2026BA}: relative distance $0.05$ with
a $20$-bit margin at $K=2^{20}$. These bounds concern each construction's
stated random ensemble. For repetition-$4$ RAA, we report Blaze's
published no-test reference bound at $K=2^{22}$ separately from our
smaller timing sizes. Its relative distance is $0.19$, with an approximately
$13$-bit margin; setup testing can improve this margin as noted in the table.
We do not transfer this bound to $K=2^{20}$ or to repetition-$3$ RAA,
nor extrapolate the BAA or EC margins to smaller lengths.

The repository documents the comparison methodology and provides the
dependency revisions, correctness checks, actual lengths, and all $5454$
timing samples for the six measured rows of
Table~\ref{tab:transposed-comparison}\repositorycite. The SPIN row is the
certificate-matched IMT series from Table~\ref{tab:finite-bch-performance},
measured under the same host, compiler version, and timing policy.
\FloatBarrier
